\section{Marktregime}
Obwohl das Gesamtsystem primär auf Methoden des maschinellen Lernens basiert, stellt die Identifikation von Marktregimen einen wichtigen regelbasierten Ansatz dar, um die Prognoseergebnisse zu kontextualisieren und die Entscheidungslogik zu verbessern.
Marktregime sind Phasen, in denen sich die statistischen Eigenschaften der Zeitreihe (z. B. Volatilität, Trend) stark ändern.
Die Erkennung dieser Regime ermöglicht es, die Vorhersagen der KI-Modelle entsprechend zu gewichten und betriebswirtschaftliche Entscheidungen gezielter zu steuern.

\subsection{Identifikation der Marktregime}
Marktphasen werden mittels eines KMeans-Clustering-Verfahrens auf Basis historischer Aluminium-Preisdaten identifiziert.
Als Eingangsgrößen dienen drei Merkmale: die EWMA-Volatilität (21-Tage-Fenster, als GARCH-Proxy), der rollende Preistrend (10-Tage-Fenster) sowie die rollende Schiefe der Preisverteilung.
Das Clustering wird ausschließlich auf den Trainingsdaten (bis einschließlich 2019) gefittet und anschließend auf den gesamten Datensatz angewendet, um eine Datenleckage in Validierungs- und Testzeitraum zu vermeiden.
Es werden drei Regime unterschieden, deren Schwerpunkte (\textit{Centroids}) in Tabelle~\ref{tab:regime_centroids} dargestellt sind.

\begin{table}[H]
\centering
\caption{Regime-Schwerpunkte (Trainingsdaten)}
\label{tab:regime_centroids}
\begin{tabular}{lrrrr}
\toprule
Regime & Volatilität & Trend & Schiefe & Beobachtungen \\
\midrule
0 & 0{,}0071 & $-$0{,}0015 & $-$0{,}40 & 1203 \\
1 & 0{,}0084 & $+$0{,}0018 & $+$0{,}83 &  901 \\
2 & 0{,}0195 & $+$0{,}0042 & $+$1{,}45 &   82 \\
\bottomrule
\end{tabular}
\end{table}

\subsubsection{Regime 0 (Stabiler Bärenmarkt)}
Dieses Regime ist durch die geringste Volatilität und einen leicht negativen Preistrend gekennzeichnet.
Die negative Schiefe deutet auf eine linksschiefe Preisverteilung hin, d.\,h.\ moderate Verluste überwiegen leichte Gewinne.
Mit rund 53\,\% aller Beobachtungen ist es das häufigste Regime.

\subsubsection{Regime 1 (Stabiler Aufwärtstrend)}
Dieses Regime weist eine vergleichbare Volatilität wie Regime~0 auf, jedoch einen positiven Preistrend und eine rechtsschiefe Verteilung.
Es macht etwa 42\,\% der Datenpunkte aus und repräsentiert geordnete Aufwärtsphasen ohne außergewöhnliche Marktverwerfungen.

\subsubsection{Regime 2 (Hochvolatile Ausnahmephase)}
Regime~2 ist durch eine rund 2,7-fach erhöhte Volatilität gegenüber Regime~0 sowie eine stark rechtsschiefe Verteilung gekennzeichnet.
Es tritt mit nur etwa 5\,\% aller Beobachtungen selten auf und korrespondiert typischerweise mit fundamentalen Marktveränderungen oder geopolitischen Schocks -- im vorliegenden Datensatz zählt die COVID-19-Pandemiephase 2020 dazu, die entsprechende Ausreißer in Validierungs- und Testergebnissen verursachen kann.

\subsection{Integration in die Modelle}
Das erkannte Regime wird nicht als regelbasierter Filter eingesetzt, sondern als numerisches kontinuierliches Merkmal dem Feature-Set der maschinellen Lernmodelle hinzugefügt.
Die Modelle lernen damit implizit, Preisprognosen in Abhängigkeit des aktuellen Marktregimes anzupassen.
Dieser Ansatz ermöglicht es insbesondere Baummodellen und neuronalen Netzen, regimespezifische Muster zu erlernen, ohne dass explizite Regelkodierungen notwendig sind.
